Ce problème se présente avec une netteté particulière pour les couplages adimensionnels.
Une coordonnée exprimée dans des unités dépend de la convention métrologique ;
une relation adimensionnelle permet d'interroger sa détermination
indépendamment de ce choix. La constante de structure fine intervient
dans des relations quantitatives susceptibles d'être confrontées à l'expérience et soulève, en même temps, la question
de la structure qui fixe le couplage électromagnétique. Le programme
historique de son explication fournit un antécédent précis pour
étudier conjointement la valeur d'une constante et l'organisation
mathématique dont elle procède.

Un épisode éclairant apparaît dans le travail de Dirac de 1931 sur les
singularités quantifiées du champ électromagnétique. Son objectif concerne
expressément l'existence d'une charge électrique minimale. Le
résultat établit une relation entre cette charge et l'intensité d'un
monopôle magnétique ; Dirac remarque toutefois que la construction ne
détermine pas à elle seule le quotient adimensionnel associé à la valeur
approximative \(137\)\cite{dirac1931}. On obtient une véritable contrainte structurelle
dont la force démonstrative n'équivaut pas à une fixation indépendante
du couplage. Ce précédent montre que la question de l'origine
des constantes appartenait à la recherche fondamentale et exige
d'identifier exactement ce que détermine chaque théorie.

L'histoire de l'électrodynamique quantique précise la différence entre
cette exigence et le calcul d'effets observables. Dans sa conférence Nobel
de 1965, Feynman a reconstitué des procédures exprimant les déplacements
de niveaux et les processus de diffusion au moyen de la masse expérimentale de
l'électron\cite{feynman1965}. La renormalisation a permis d'obtenir
des résultats finis et confrontables à l'expérience tout en conservant cette référence. L'efficacité
de ces prédictions constitue un acquis théorique effectif ; leur dépendance
à l'égard d'un paramètre expérimental identifie, en même temps, la place logique
où une détermination supplémentaire pourrait élargir l'explication. La
question constitutive s'ajoute à la capacité prédictive et étudie les
conditions qui fixent ses paramètres.

Le rapport conjoint LEP--SLD de 2006 fournit une manifestation expérimentale
de cette architecture\cite{lep2006}. Ses auteurs distinguent les paramètres
libres du modèle des relations que celui-ci impose entre observables.
Par les corrections radiatives, les données de la résonance Z permettaient
d'inférer indirectement les masses du quark top et du boson W et de les confronter
aux mesures directes. Les prédictions entrelacées et la détermination
expérimentale des paramètres formaient une même procédure de mise à l'épreuve.
Le problème pertinent consiste donc à identifier les dépendances qu'une
construction détermine intérieurement et celles dont la spécification continue
de nécessiter une information empirique.

Dans ce contexte, on entendra par \emph{détermination épistémologique}
la procédure qui établit une valeur par la relation entre théorie,
observation et inférence ; par \emph{détermination constitutive}, la
déduction de cette valeur à partir des relations définissant l'objet
considéré. Cette seconde expression désigne ici une exigence mathématique :
expliciter le domaine, les opérateurs, les normalisations et les
compatibilités sélectionnant le résultat. Sa pertinence physique
procède de la confrontation ultérieure aux grandeurs observées.
Le changement de perspective concerne l'ordre de l'explication : la valeur est
examinée comme conséquence d'une structure déterminée. L'aspiration
ontologique consiste à relier cette structure à la constitution de
l'objet physique ; la preuve mathématique et sa confrontation expérimentale établissent
la portée avec laquelle une telle relation peut être soutenue.

